\section{AnyBlox}

A future-proof data access solution must satisfy two conceptual requirements simultaneously. First, it must introduce an abstraction layer that decouples format internals from processing systems and system internals from decoders, thereby reducing the $N \times M$ integration burden to $N + M$. Second, this abstraction must not preclude direct file access, since routing data through an intermediary process introduces transfer overhead that can dominate total runtime. These two goals are in tension: an abstraction layer must be located somewhere, and every candidate location either breaks direct access or reintroduces format-specific coupling~\cite{gienieczkoAnyBloxFrameworkSelfDecoding2025}.

Existing approaches each resolve this tension in favor of one side. Native integration preserves direct file access and achieves the best performance through tight coupling with host internals, but requires a separate decoder implementation per format per system, thus reproducing the $N \times M$ problem. Extension mechanisms reduce the per-format implementation effort, but each host exposes its own extension API, so a decoder written for system A cannot be reused in system B without a near-complete rewrite. Isolated extensions, such as containerized user-defined functions \cite{saurContainerizedExecutionUDFs2022a} or remote cloud user-defined code execution solutions provide both isolation and portability, but introduce a data transfer bottleneck between the host and the isolated decoder process~\cite{gienieczkoAnyBloxFrameworkSelfDecoding2025}.

Given these constraints, Gienieczko et al. (2025) argue that there is only one location for the decoder that satisfies both requirements: inside the data file itself. Packaging the decoder with the data eliminates the transfer bottleneck, since decoder and data reside in the same address space, while the abstraction is preserved because the host need not implement any format-specific logic. This is the core idea behind self-decoding data and the design premise of AnyBlox.

\subsection{WebAssembly as the Decoder Runtime}

A self-decoding dataset requires a decoder runtime that satisfies four criteria: \textbf{Portability}, \textbf{Security}, \textbf{Performance} and \textbf{Extensibility}.

Gienieczko et al. (2025) identify WebAssembly (Wasm) as a runtime that satisfies all four simultaneously. Wasm is a portable low-level bytecode and virtual machine specification, designed as an abstraction over modern hardware rather than a commitment to any specific programming model, making it language-, hardware-, and platform-independent~\cite{haasBringingWebSpeed2017}.

\subsubsection{Portability}

A decoder implemented in Wasm runs on a wide variety of host architectures (e.g. browsers, mobile devices, embedded, and server workstations)~\cite{gienieczkoAnyBloxFrameworkSelfDecoding2025}. For AnyBlox, this means a format author writes and compiles a decoder once, and any host can execute it via \texttt{libanyblox}, independent of the underlying hardware or operating system.

\subsubsection{Security}

The WebAssembly specification provides machine-verifiable memory safety and isolation guarantees~\cite{wattMechanisingVerifyingWebAssembly2018}. AnyBlox reinforces this by being implemented predominantly in safe Rust, guaranteeing memory and thread safety. All unsafe code is confined to a single module handling Wasm memory maps, keeping it easily auditable~\cite{gienieczkoAnyBloxFrameworkSelfDecoding2025}.

\subsubsection{Performance}

Wasm is JIT-compiled to the host's native instruction set at runtime, making near-native performance attainable. Additionally, AnyBlox avoids unnecessary data copying through a custom memory management scheme. Gienieczko et al. (2025) note that AnyBlox can passively benefit from future improvements in Wasm compiler technology, since all JIT and sandboxing details are encapsulated in \texttt{libanyblox} and not exposed in the public API~\cite{gienieczkoAnyBloxFrameworkSelfDecoding2025}.

\subsubsection{Extensibility}

Since most general-purpose programming languages provide a Wasm toolchain, format authors can implement decoders in their language of choice without restrictions. The main practical obstacle is SIMD instructions: converting platform-specific SIMD intrinsics to Wasm's instruction set requires manual developer effort. However, since Wasm defines SIMD at the bytecode level, the resulting modules remain fully portable across architectures~\cite{gienieczkoAnyBloxFrameworkSelfDecoding2025}.